\chapter{単位・計測}
\section{背景：量の意味と数値を共有するための単位}
観測量を数値で比較するには、同じ物理量を同じ尺度で表す必要がある。単位の定義は尺度の基準を定め、実現は実験でその尺度へ接続する方法である。定義・物理法則・測定手順を同じものとして扱わない。
SIは長さm、質量kg、時間s、電流A、熱力学温度K、物質量mol、光度cdの7基本単位を持つ。平成23・27年度の原典の導入はcdを欠く6単位の列挙であり、光度を含む全物理量をその6単位だけで組み立てられるという説明は保持できない。誤った選択肢の判定と、この導入の疑義を別に記録する\cite{si_brochure9}。
本章は対象11問の原典指数・条件を確認した初稿である。正答を導けても、2導入と改定後の三重点の厳密さ1問には疑義を残す。Astra本文査読前である。

\section{量方程式と数値方程式}\label{theory:units:quantity}
量は「数値×単位」である。量方程式 \(F=ma\) は単位の選択に依存しないが、数値だけの式は単位を揃える必要がある。
量の次元を \(M,L,T,I\) 等で表せば、力は \(MLT^{-2}\)、エネルギーは \(ML^2T^{-2}\) である。加算する量は同次元でなければならず、指数・対数の引数は無次元である。
次元一致は必要条件であって十分条件ではない。同じ次元を持つ仕事とトルク、周波数と放射能でも物理的意味は異なる。定数・係数・方向や対象を次元だけで確定しない。

\section{力学量：力から圧力・仕事・仕事率へ}\label{theory:units:derived}
運動方程式、面積への力、仕事、単位時間の仕事から順に
\[
 [F]=\mathrm{kg\,m\,s^{-2}}=\mathrm N,\qquad
 [p]=[F]/[S]=\mathrm{kg\,m^{-1}s^{-2}}=\mathrm{Pa},
\]
\[
 [E]=[F][\ell]=\mathrm{kg\,m^2s^{-2}}=\mathrm J,\qquad
 [P]=[E]/[t]=\mathrm{kg\,m^2s^{-3}}=\mathrm W
\]
となる。仕事は一般には \(\int\boldsymbol F\cdot d\boldsymbol r\) であるため、単純積は一定力とその方向の移動距離の場合を表す。次元の導出には同じ長さ因子を使う。
平成27年度のWの長さ指数1は、仕事を作る移動距離を一つ欠く。Paは面積で割るのでmの指数が \(-1\) になる。力の単位をそのまま圧力へ使わない\cite{si_nist811}。

\section{電気量：電荷・電圧・抵抗・磁束密度}\label{theory:units:electrical}
一定電流なら \(Q=It\)、一般には \(Q=\int I\,dt\) でありC=Asである。電位差は電荷当たりの仕事なので
\[
 [V]=[\text{仕事}]/[Q]=\mathrm{kg\,m^2s^{-3}A^{-1}},\qquad
 [R]=[V]/[I]=\mathrm{kg\,m^2s^{-3}A^{-2}}.
\]
原典のVのm指数1やΩのA指数 \(-1\) はこの接続を欠く。電場は力/電荷でN/C、電位差はJ/Cであり、両者は長さ一つ違う。
一様磁場に垂直な直線電流の磁気力 \(F=I\ell B\) から
\[
 [B]=\mathrm{N/(A\,m)}=\mathrm{kg\,s^{-2}A^{-1}}=\mathrm T
\]
である。テスラへ不要なmを残したりAを二回で割ったりしない。単位はオームの法則の適用材料・一定温度等の条件を保証するものではない\cite{si_nist811}。

\section{周波数・角度・放射能の意味}\label{theory:units:angle}\label{theory:units:activity}
周期運動の周波数は \(f=1/T\) でHz=\(\mathrm{s^{-1}}\)、平面角は弧長/半径でradの次元は1である。角の意味を残すためradを用い、角周波数 \(\omega\) と周期周波数 \(f\) を \(\omega=2\pi f\) で区別する。
放射能は核壊変数の時間率でBq=\(\mathrm{s^{-1}}\) である。HzとBqは同次元だが、周期数と壊変数の意味は異なる。線量GyやSvに変換するには壊変ごとのエネルギー、吸収先・質量等の追加情報が必要である。
仕事とトルクにも同じ次元が現れる。数値・次元が一致しても、何を測る量かを読み落とさない\cite{si_brochure9}。

\section{定義の改定：原器から固定した定数へ}\label{theory:units:history}
単位を人工物へ結び付けると、長期変化や各国への伝達に限界がある。再現可能な自然定数と原子の周波数へ結び付けることで、各地の実験から一つの尺度へ接続できる。
\begin{tcolorbox}[title=出題時の定義を固定する]
平成28・29年度は2019年5月20日より前である。kgは国際キログラム原器、Kは規定水の三重点、molは0.012kgの炭素12中の原子数と等しい数の指定要素粒子に基づいた。これらの出題時の適切肢を、現在の定義だけで誤りへ変えない。
\end{tcolorbox}
令和元年度再試験は問題幹が2019年5月の改定後を明記する。原器によるkg定義はそこで不適切となる。改定によって既存の測定値が一斉に別の数値へ跳ぶように設計されたわけではない。連続性を保ちつつ、以前固定されていた物質の量に測定不確かさが移る\cite{si_cgpm2018,si_nist811}。

\section{7基本単位と定義定数}\label{theory:units:definition}
現行SIの定義定数との接続は次の通りである。定数と単位が互いに孤立した一対一の定義ではなく、既に固定した時間・長さ等を組み合わせる点に注意する。
\begin{center}
\begin{tabular}{lll}
量・基本単位 & 主な定義定数 & 固定する数値\\\hline
時間 s & Cs133超微細遷移周波数 & \(9192631770\,\mathrm{Hz}\)\\
長さ m & 真空光速 \(c\) & \(299792458\,\mathrm{m/s}\)\\
質量 kg & プランク定数 \(h\) & \(6.62607015\times10^{-34}\,\mathrm{J\,s}\)\\
電流 A & 素電荷 \(e\) & \(1.602176634\times10^{-19}\,\mathrm C\)\\
熱力学温度 K & ボルツマン定数 \(k\) & \(1.380649\times10^{-23}\,\mathrm{J/K}\)\\
物質量 mol & アボガドロ定数 \(N_{\rm A}\) & \(6.02214076\times10^{23}\,\mathrm{mol^{-1}}\)\\
光度 cd & 視感効果度 \(K_{\rm cd}\) & \(683\,\mathrm{lm/W}\)（540THz）\\
\end{tabular}
\end{center}
たとえば \(h\) の次元は \(\mathrm{kg\,m^2/s}\) であるため、mとsの定義を踏まえてkgの尺度を得る。重力定数 \(G\) をkgの定義に固定しているわけではない。kgを天秤の重さ測定へ実現する際に重力加速度を使うこととは区別する。
電流には \(e\) と秒、温度には \(k\) とエネルギー尺度、物質量には指定粒子数を使う。今回の6定義問題は、改定の時点又はこの対応の取り違えを問う\cite{si_cgpm2018}。

\section{水の三重点・凝固点・Celsius・温度尺度}\label{theory:units:temperature}
三重点は固体・液体・気体の三相が平衡にある温度と圧力の点である。通常圧力での凝固点とは違う。Celsius温度の数値は \(t/{}^\circ\mathrm C=T/\mathrm K-273.15\) である。
旧K定義では規定同位体組成の水の三重点は厳密に273.16K、したがって0.01℃であった。水組成・圧力の条件を省き、すべての水の凝固がその値になるとしない\cite{si_kelvin2007}。
\begin{tcolorbox}[title=改定後の三重点は熱力学温度の厳密定義値ではない]
2019年後のKは \(k\) の固定で定義する。BIPM実現文書は三重点の最良推定273.16K、標準不確かさ0.1mKを示す。実用温度尺度ITS-90の \(T_{90}\) の固定点と、熱力学温度 \(T\) の厳密定義を区別する。令和元年度再試験の0.01℃は提示桁の近似として理解できるが、厳密な定義値という読みに疑義を残す。
\end{tcolorbox}
「値がほぼ同じ」と「定義が同じ」を混同しない\cite{si_kelvin_realization}。

\section{光度：黒体基準と単色放射の尺度}\label{theory:units:photometry}
放射エネルギーの尺度だけでは人の視覚に対する明るさを同じ数値で表せない。周波数による視感度を重み付けした測光量が必要になる。光度cdは方向当たりの光束であり、エネルギー流のWや面積当たりの照度lxとは異なる。
現在は周波数 \(540\times10^{12}\,\mathrm{Hz}\) の単色放射に対する \(K_{\rm cd}=683\,\mathrm{lm/W}\) を固定する。その放射の放射強度が \(1/683\,\mathrm{W/sr}\) の方向の光度が1cdに対応する。
黒体を使った歴史的定義と現在の定義、黒体光源を使う実験上の実現は別である。令和3年度の「黒体で定義」は不適切、令和7年度の「単色放射で定義」はこの固定周波数・視感度の省略として適切である\cite{si_brochure9}。

\section{換算：面積指数・力と質量・動粘度}\label{theory:units:conversion}
単位を換えるときは量を変えず、長さ換算を指数の回数だけ適用する。\(1\,\mathrm{cm^2}=10^{-4}\,\mathrm{m^2}\) であり \(10^{-2}\) ではない。
kgfは標準重力加速度 \(g_{\rm n}=9.80665\,\mathrm{m/s^2}\) により \(1\,\mathrm{kgf}=9.80665\,\mathrm N\) とした力の単位であり、kgは質量である。したがって
\[
 1\,\mathrm{kgf/cm^2}=\frac{9.80665}{10^{-4}}\,\mathrm{N/m^2}
 =9.80665\times10^4\,\mathrm{Pa}.
\]
標準atmは101325Pa、barは \(10^5\) Pa、通常のmmHgは約133.322Pa、psiは約6894.757Paである。液柱の実際の圧力 \(\rho g h\) は流体の温度や重力条件に依存する。通常の換算表の係数と実測条件を区別する。
動粘度 \(\nu=\mu/\rho\) の次元は \(L^2/T\)。Stは \(\mathrm{cm^2/s}\) なので \(10^{-4}\,\mathrm{m^2/s}\) となる。非SI単位であっても、そのSIへの換算が正しい場合がある\cite{si_nist811}。

\section{今回の問われ方と学習順序}
主概念は組立単位・換算5問と定義・改定6問である。まず物理量の量方程式からm・s・Aの指数を導き、次に定義の基準日を置き、最後に非SI係数を照合する。単位を丸暗記すると、原器と三重点の改定や電場と電位の違いを取りこぼす。
この集計は本章11問の確認結果であり将来の出題保証ではない。疑義の3問も含めて全候補を追跡する。

\section{過去問と全候補の検討}
\subsection{平成23年度 問6}
\label{q:h23-reg-applied-06}

\noindent\textbf{原典：} sources/H23/17\_applied\_science.pdf、PDF物理3頁。

\textbf{公式正答：} 3（single）。

\par\textit{原典本文の校正と独自解説は未作成。分類のみの登録項目。}

\subsection{平成25年度 問4}
\label{q:h25-reg-applied-04}

\noindent\textbf{原典：} sources/H25/17\_applied\_science.pdf、PDF物理3頁。

\textbf{公式正答：} 5（single）。

\par\textit{原典本文の校正と独自解説は未作成。分類のみの登録項目。}

\subsection{平成27年度 問4}
\label{q:h27-reg-applied-04}

\noindent\textbf{原典：} sources/H27/17\_applied\_science.pdf、PDF物理2頁。

\textbf{公式正答：} 4（single）。

\par\textit{原典本文の校正と独自解説は未作成。分類のみの登録項目。}

\subsection{平成28年度 問6}
\label{q:h28-reg-applied-06}

\noindent\textbf{原典：} sources/H28/17\_applied\_science.pdf、PDF物理2頁。

\textbf{公式正答：} 4（single）。

\par\textit{原典本文の校正と独自解説は未作成。分類のみの登録項目。}

\subsection{平成29年度 問5}
\label{q:h29-reg-applied-05}

\noindent\textbf{原典：} sources/H29/17\_applied\_science.pdf、PDF物理2頁。

\textbf{公式正答：} 4（single）。

\par\textit{原典本文の校正と独自解説は未作成。分類のみの登録項目。}

\subsection{令和元年度再試験 問1}
\label{q:r01-ret-applied-01}

\noindent\textbf{原典：} sources/R01\_retest/17\_applied\_science.pdf、PDF物理1頁。

\textbf{公式正答：} 1（single）。

\par\textit{原典本文の校正と独自解説は未作成。分類のみの登録項目。}

\subsection{令和2年度 問9}
\label{q:r02-reg-applied-09}

\noindent\textbf{原典：} sources/R02/17\_applied\_science.pdf、PDF物理3頁。

\textbf{公式正答：} 2（single）。

\par\textit{原典本文の校正と独自解説は未作成。分類のみの登録項目。}

\subsection{令和3年度 問9}
\label{q:r03-reg-applied-09}

\noindent\textbf{原典：} sources/R03/17\_applied\_science.pdf、PDF物理4頁。

\textbf{公式正答：} 1（single）。

\par\textit{原典本文の校正と独自解説は未作成。分類のみの登録項目。}

\subsection{令和5年度 問13}
\label{q:r05-reg-applied-13}

\noindent\textbf{原典：} sources/R05/17\_applied\_science.pdf、PDF物理5頁。

\textbf{公式正答：} 5（single）。

\par\textit{原典本文の校正と独自解説は未作成。分類のみの登録項目。}

\subsection{令和6年度 問16}
\label{q:r06-reg-applied-16}

\noindent\textbf{原典：} sources/R06/17\_applied\_science.pdf、PDF物理5頁。

\textbf{公式正答：} 3（single）。

\par\textit{原典本文の校正と独自解説は未作成。分類のみの登録項目。}

\subsection{令和7年度 問9}
\label{q:r07-reg-applied-09}

\noindent\textbf{原典：} sources/R07/17\_applied\_science.pdf、PDF物理3頁。

\textbf{公式正答：} 3（single）。

\par\textit{原典本文の校正と独自解説は未作成。分類のみの登録項目。}

